
\documentclass[12pt,a4paper]{article}

% =====================================================
% PAKETI I PODESAVANJA
% =====================================================

\usepackage[T1]{fontenc}
\usepackage[utf8]{inputenc}
\usepackage[serbian]{babel}
\usepackage{lmodern}

\usepackage[a4paper,margin=2.5cm]{geometry}
\usepackage{graphicx}
\usepackage{float}
\usepackage{enumitem}
\usepackage{hyperref}
\usepackage{setspace}

\onehalfspacing
\setlist{itemsep=2pt,topsep=4pt}

\hypersetup{
    colorlinks=true,
    linkcolor=black,
    urlcolor=blue,
    citecolor=black
}

\begin{document}

% =====================================================
% UVOD
% =====================================================

\section*{Uvod}

Dokument opisuje tri slučaja upotrebe informacionog sistema za
rent-a-car agenciju: upravljanje korisnicima, upravljanje iznajmljivanjima
i preuzimanje i vraćanje vozila. Za svaki slučaj navedeni su akteri,
preduslovi, osnovni i alternativni tokovi, kao i postuslov.


% =====================================================
% 1. SLUCAJ UPOTREBE
% =====================================================

\section*{1. Slučaj upotrebe: Upravljanje korisnicima}

\subsection*{Kratak opis}

Zaposleni koristi aplikaciju za upravljanje podacima o korisnicima. Može da
kreira novog korisnika, pregleda i pretražuje postojeće korisnike, menja njihove
podatke i briše, odnosno deaktivira zapis kada za to postoje uslovi.

\subsection*{Akteri}

\begin{itemize}
    \item Zaposleni -- radnik rent-a-car agencije koji koristi aplikaciju.
\end{itemize}

\subsection*{Preduslov}

\begin{itemize}
    \item Zaposleni mora imati kreiran nalog.
    \item Zaposleni mora biti uspešno prijavljen u sistem.
\end{itemize}

\subsection*{Osnovni tok}

\begin{enumerate}
    \item Zaposleni otvara sekciju „Upravljanje korisnicima“.
    \item Sistem prikazuje listu evidentiranih korisnika i dostupne opcije.
    \item Zaposleni unosi kriterijum za pretragu ili bira opciju za dodavanje korisnika.
    \item Sistem prikazuje rezultate pretrage ili formu za unos novog korisnika.
    \item Zaposleni bira postojećeg korisnika radi pregleda, izmene ili brisanja,
    odnosno unosi podatke novog korisnika.
    \item Zaposleni unosi ili menja tražene podatke i potvrđuje akciju.
    \item Ako je izabrano brisanje, sistem traži dodatnu potvrdu.
    \item Sistem proverava ispravnost podataka i uslove za izmenu ili brisanje.
    \item Sistem čuva promene u bazi podataka i prikazuje rezultat akcije.
\end{enumerate}


\subsection*{Alternativni tokovi}

\begin{description}[style=nextline]

    \item[A1: Korisnik nije pronađen.]
    Sistem prikazuje poruku da nema korisnika koji odgovaraju kriterijumu pretrage.
    Zaposleni se vraća na korak 3 i unosi novi kriterijum ili bira dodavanje korisnika.

    \item[A2: Neispravni ili nepotpuni podaci.]
    Prilikom provere (korak 8) sistem prikazuje grešku i označava polja koja treba ispraviti ili dopuniti. Zaposleni se vraća na korak 5 i ponavlja unos.

    \item[A3: Korisnik već postoji.]
    Prilikom provere (korak 8) sistem utvrđuje da korisnik postoji, prikazuje upozorenje i ne kreira duplikat. Zaposleni se vraća na korak 5 da ispravi podatke, a može i da odustane (A5).

    \item[A4: Korisnik ima aktivno iznajmljivanje.]
    Prilikom provere brisanja (korak 8) sistem ne dozvoljava brisanje dok postoji aktivno iznajmljivanje i prikazuje poruku. Zaposleni se vraća na korak 5 i bira drugu akciju ili drugog korisnika, a može i da odustane (A5).

    \item[A5: Odustajanje od akcije.]
    Zaposleni ne potvrđuje akciju (korak 6 ili 7). Sistem ne čuva izmene, baza ostaje nepromenjena i tok se završava. Zaposleni ostaje na pregledu korisnika.

\end{description}


\subsection*{Postuslov}

Podaci o korisniku su kreirani, pregledani, izmenjeni ili deaktivirani.
Ako je akcija otkazana ili nije uspešna, baza ostaje nepromenjena.


\begin{figure}[H]
    \centering
    \includegraphics[width=0.9\textwidth]{upravljanje korisnicima.jpg}
    \caption{Dijagram slučaja upotrebe za upravljanje korisnicima}
    \label{fig:uml-korisnici}
\end{figure}

\begin{figure}[H]
    \centering
    \includegraphics[width=0.9\textwidth]{dijagram_aktivnosti_korisnici.png}
    \caption{Dijagram aktivnosti korisnika}
    \label{fig:aktivnosti-korisnici}
\end{figure}

% 2. SLUCAJ UPOTREBE
% =====================================================

\section*{2. Slučaj upotrebe: Upravljanje iznajmljivanjima}

\subsection*{Kratak opis}

Zaposleni upravlja iznajmljivanjima vozila. Može da pretraži postojeća
iznajmljivanja, kreira novo iznajmljivanje za korisnika (uz proveru
dostupnosti vozila i obračun cene), izmeni njegove podatke ili ga
otkaže kada je to dozvoljeno.

\subsection*{Akteri}

\begin{itemize}
    \item Zaposleni -- radnik rent-a-car agencije koji evidentira i ažurira iznajmljivanja.
\end{itemize}

\subsection*{Preduslov}

\begin{itemize}
    \item Zaposleni mora biti uspešno prijavljen u sistem.
    \item Korisnik i vozilo moraju biti evidentirani u sistemu.
\end{itemize}

\subsection*{Osnovni tok (kreiranje novog iznajmljivanja)}

\begin{enumerate}
    \item Zaposleni otvara sekciju „Upravljanje iznajmljivanjima“.
    \item Sistem prikazuje listu iznajmljivanja i dostupne opcije.
    \item Zaposleni bira korisnika, vozilo i unosi datum i vreme preuzimanja i vraćanja.
    \item Sistem proverava ispravnost perioda i podataka.
    \item Sistem proverava da li je vozilo slobodno u izabranom periodu.
    \item Ako je vozilo dostupno, sistem obračunava cenu i prikazuje je zaposlenom.
    \item Zaposleni proverava podatke i potvrđuje iznajmljivanje.
    \item Sistem čuva iznajmljivanje i dodeljuje mu status „Kreirano“.
\end{enumerate}

\subsection*{Alternativni tokovi}

\begin{description}[style=nextline]

    \item[A1: Vozilo nije dostupno.]
    Sistem prikazuje obaveštenje o zauzetosti vozila.
    Zaposleni se vraća na korak 3 i bira drugo vozilo ili menja period iznajmljivanja.

    \item[A2: Neispravan period.]
    Sistem prikazuje poruku o grešci da vreme vraćanja mora biti nakon vremena preuzimanja.
    Zaposleni se vraća na korak 3 i ispravlja period.

    \item[A3: Neispravni ili nepotpuni podaci.]
    Sistem prikazuje poruku o grešci i označava podatke koje je potrebno ispraviti ili dopuniti.
    Zaposleni se vraća na korak 3 i ponavlja unos.

    \item[A4: Izmena iznajmljivanja.]

    Zaposleni pretražuje i bira postojeće iznajmljivanje, a sistem proverava njegov status.
    Ako je akcija dozvoljena, zaposleni bira izmenu i menja dozvoljene podatke.
    Sistem ponovo proverava dostupnost vozila i obračunava cenu.
    Zaposleni potvrđuje izmenu.
    Sistem čuva izmene i prikazuje rezultat.

    \item[A5: Otkazivanje iznajmljivanja.]

    Zaposleni pretražuje i bira postojeće iznajmljivanje, a sistem proverava njegov status.
    Ako je akcija dozvoljena, zaposleni pokreće otkazivanje i potvrđuje ga.
    Sistem menja status iznajmljivanja u „Otkazano“ i prikazuje rezultat.

    \item[A6: Nedozvoljena izmena ili otkazivanje.]
    Ako je iznajmljivanje završeno ili već otkazano, sistem odbija akciju.
    Tok se završava, a podaci ostaju nepromenjeni.

    \item[A7: Iznajmljivanje nije pronađeno.]
    Prilikom pretrage (A4, A5) sistem prikazuje poruku da nema rezultata.
    Zaposleni se vraća na izbor iznajmljivanja i bira drugo.

\end{description}

\subsection*{Postuslov}

Novo iznajmljivanje je sačuvano sa statusom „Kreirano“, postojeće je
ažurirano ili je dobilo status „Otkazano“. Ako akcija nije uspešno
završena, podaci ostaju nepromenjeni.



\begin{figure}[H]
    \centering
    \includegraphics[width=0.9\textwidth]{upravljanje iznajmljivanjima.jpg}
    \caption{Dijagram slučajeva upravljanje iznajmljivanjima}
    \label{fig:aktivnosti-kreiranje-iznajmljivanja}
\end{figure}

\begin{figure}[H]
    \centering
    \includegraphics[width=0.9\textwidth]{dijagram aktivnosti iznajmljivanja.png}
    \caption{Dijagram aktivnosti iznajmljivanja}
    \label{fig:aktivnosti-izmena-otkazivanje}
\end{figure}

\begin{figure}[H]
    \centering
    \includegraphics[width=0.9\textwidth]{dijagram aktivnosti izemena i otkazivanje (1).png}
    \caption{Dijagram aktivnosti izmene i otkazivanja}
    \label{fig:stanja-iznajmljivanja}
\end{figure}

\begin{figure}[H]
    \centering
    \includegraphics[width=0.9\textwidth]{dijagram stanja iznajmljivanja.png}
    \caption{Dijagram stanja iznajmljivanja}
    \label{fig:stanja-iznajmljivanja}
\end{figure}



% =====================================================
% 3. SLUCAJ UPOTREBE
% =====================================================

\section*{3. Slučaj upotrebe: Preuzimanje i vraćanje vozila}

\subsection*{Kratak opis}

Zaposleni evidentira predaju vozila korisniku i njegovo vraćanje nakon završetka
iznajmljivanja. Prilikom preuzimanja i vraćanja proverava stanje vozila,
evidentira eventualna oštećenja i proverava da li je vozilo vraćeno u
ugovorenom roku. Ako postoji kašnjenje, sistem obračunava penale prema
definisanim pravilima.

\subsection*{Akteri}

\begin{itemize}
    \item Zaposleni -- radnik koji predaje vozilo korisniku i evidentira njegovo vraćanje.
\end{itemize}

\subsection*{Preduslov}

\begin{itemize}
    \item Zaposleni mora biti uspešno prijavljen u sistem.
    \item Za preuzimanje mora postojati odgovarajuće iznajmljivanje u statusu „Kreirano“.
    \item Za vraćanje mora postojati aktivno iznajmljivanje.
    \item Vozilo mora biti evidentirano u sistemu.
\end{itemize}

\subsection*{Osnovni tok}

\begin{enumerate}
    \item Zaposleni otvara sekciju „Preuzimanje i vraćanje vozila“.
    \item Zaposleni pronalazi iznajmljivanje, a sistem prikazuje podatke o korisniku,
    vozilu i periodu iznajmljivanja.
    \item Prilikom preuzimanja zaposleni pregleda vozilo, evidentira početno stanje
    i predaje vozilo korisniku.
    \item Sistem čuva podatke o pregledu i menja status iznajmljivanja u „Aktivno“.
    \item Prilikom vraćanja zaposleni pronalazi aktivno iznajmljivanje i evidentira
    stvarni datum i vreme vraćanja, kao i stanje vozila nakon korišćenja.
    \item Ako je uočeno oštećenje, zaposleni ga evidentira.
    \item Sistem upoređuje stvarno vreme vraćanja sa ugovorenim rokom.
    \item Ako postoji kašnjenje, sistem ga evidentira i obračunava penale prema pravilima.
    \item Sistem čuva podatke o vraćanju, oštećenjima i penalima.
    \item Sistem menja status iznajmljivanja u „Završeno“ i ažurira dostupnost vozila.
\end{enumerate}

\subsection*{Alternativni tokovi}

\begin{description}[style=nextline]

    \item[A1: Ne postoji odgovarajuće iznajmljivanje.]
    Nakon pretrage (korak 2) sistem obaveštava zaposlenog da za izabranog korisnika
    ili vozilo ne postoji odgovarajuće iznajmljivanje. Tok se završava, a podaci
    ostaju nepromenjeni.

    \item[A2: Uočeno je oštećenje.]
    U koraku 6 zaposleni opisuje oštećenje i unosi potrebne podatke. Sistem čuva
    zapis povezan sa vozilom i iznajmljivanjem zajedno sa podacima o vraćanju
    (korak 9), a tok se nastavlja korakom 7.

    \item[A3: Vozilo je vraćeno nakon roka.]
    U koraku 8 sistem utvrđuje trajanje kašnjenja, obračunava penale i prikazuje
    iznos zaposlenom. Tok se nastavlja korakom 9.

    \item[A4: Nije moguće obračunati penale.]
    Ako prilikom obračuna (korak 8) tarifa nije definisana, sistem obaveštava
    zaposlenog da je potrebna dopuna podataka ili ručni postupak. Nakon toga se
    tok nastavlja korakom 9.

    \item[A5: Nema oštećenja i nema kašnjenja.]
    Nema oštećenja (korak 6) i nema kašnjenja (korak 8), pa se preskaču
    obračun penala i evidentiranje oštećenja. Sistem evidentira uredno vraćanje
    bez dodatnih troškova, a tok se nastavlja korakom 9.

\end{description}

\subsection*{Postuslov}

Preuzimanje ili vraćanje vozila je evidentirano. Nakon preuzimanja iznajmljivanje
dobija status „Aktivno“, a nakon vraćanja status „Završeno“. Vozilo se označava
kao dostupno ako nema drugih razloga koji to sprečavaju. Podaci o oštećenjima i
penalima čuvaju se uz odgovarajuće iznajmljivanje.


\begin{figure}[H]
    \centering
    \includegraphics[width=0.9\textwidth]{preuzimanje i vracanje vozila.jpg}
    \caption{Dijagram slučaja za preuzimanje i vraćanje vozila}
    \label{fig:aktivnosti-kreiranje-iznajmljivanja}
\end{figure}

\begin{figure}[H]
    \centering
    \includegraphics[width=0.9\textwidth]{dijagram aktivnosti vracanje vozila.png}
    \caption{Dijagram aktivnosti za preuzimanje i vraćanje vozila}
    \label{fig:aktivnosti-izmena-otkazivanje}
\end{figure}


\end{document}